% Proved after four complete actual-file/source audits. Explicit geometric-period diagnostic on one fixed complete elliptic cone; no representation-count or simple-quotient conclusion.
\subsection{Unbounded unitary geometric periods on one elliptic cone}

\begin{lemma}[Proved: explicit tame Heisenberg covers of a fixed cone]
\label{D:elliptic-cone-Heisenberg-covers}
Let $k$ be algebraically closed of characteristic $p>0$, let
$(E,O)$ be an elliptic curve over $k$, and put
$N=\mathcal O_E(-3O)$. Let $R$ be the completion at the vertex
of the affine cone for the plane embedding given by
$N^{-1}=\mathcal O_E(3O)$, and put
$U=\operatorname{Spec}R\setminus\{\mathfrak m_R\}$.
This is one fixed complete normal surface germ. Its resolution
has the single smooth exceptional curve $E$, with $E^2=-3$.

For every odd $m$ with $(m,3p)=1$ there is a connected finite
\'etale Galois cover $U_m\to U$ whose group has the presentation
\[
 G_m=\langle x,y,z\mid x^m=y^m=z^m=1,\quad
 z\text{ central},\quad [x,y]=z^3\rangle,
 \qquad |G_m|=m^3.                                      \tag{D.EC.1}
\]
On the corresponding regular resolution cover, the exceptional
curve is again $E$; its map to the original exceptional curve is
$[m]:E\to E$. The generic inertia group is $\langle z\rangle$
of order $m$, and the residue extension has group $E[m]$ of
order $m^2$. After choosing a primitive $m$th root $\zeta\in k$,
the generator $z$ may be chosen to act on a normal parameter
by $z(s)=\zeta s$.
\end{lemma}

\begin{proof}
\emph{Known inputs.}
We use the theta-group construction and its nondegenerate
commutator pairing for an ample line bundle of separable type
from \citet[\S1, pp.~288--295, especially Theorem~1 and its proof]
{D_Mumford1966ThetaCone}.
The multiplication and symmetric-pullback formulas are
\citet[Lemmas~39.9.7--39.9.9 and Proposition~39.9.11]
{D_StacksEllipticThetaCone}.
For an elliptic curve, the identification of $E$ with
$\operatorname{Pic}^0(E)$ through $P\mapsto\mathcal O_E(P-O)$
also follows from the genus-one case of Lemma~44.6.7 of that
reference: the degree-one Abel map has one-point projective-space
fibres. These are cited results; no prime-order assumption is
introduced below.

Write $Y_0=\operatorname{Tot}(N)$ and let $E\subset Y_0$ be its
zero section. The usual blowup charts of a homogeneous plane
cone identify $Y_0$ with the blowup of its vertex. To make the
geometry explicit, if $F(X_0,X_1,X_2)$ is the smooth cubic
equation of $E$, the chart with direction $X_i\ne0$ has equation
$F(1,u,v)=0$ and the free fibre coordinate $t=X_i$.
The other charts glue by the transition functions of
$\mathcal O_E(-1)=N$.
The hypersurface cone is a domain and Cohen--Macaulay; its
puncture is smooth by the same charts. It is therefore normal
by the codimension-one normality criterion. Its excellent
local ring has normal, hence integral, completion
\citep[Lemma~15.53.6, Tag~0C23]{D_StacksEllipticThetaCone}.
Completion preserves the regular proper resolution and its
closed fibre, by Lemmas~54.11.1--54.11.2 of the same reference.
Thus the asserted fixed complete germ and self-intersection
are as stated.

Put $M=N^{\otimes m}$. Since $N$ is symmetric,
\[
        [m]^*N\simeq N^{\otimes m^2}=M^{\otimes m}.
                                                               \tag{D.EC.2}
\]
Fix such an isomorphism $\psi$. Fibrewise $m$th power followed
by $\psi$ defines
\[
 f_m:\operatorname{Tot}(M)\longrightarrow\operatorname{Tot}(N),
 \qquad (P,w)\longmapsto([m]P,\psi_P(w^{\otimes m})).
                                                               \tag{D.EC.3}
\]
It factors as the finite locally free $m$th-power map to
$\operatorname{Tot}(M^{\otimes m})$, then the base change of
the finite \'etale isogeny $[m]$, of degree $m^2$.
Consequently $f_m$ is finite locally free of degree $m^3$,
and is \'etale off the zero section. Its source is smooth
and integral, and its inverse image of $E$ has reduced
support the single curve $E$.

For $u\in E[m]$, translation by $u$ preserves $M$:
under $E\simeq\operatorname{Pic}^0(E)$ the difference
$\tau_u^*M\otimes M^{-1}$ is the point $3mu=0$, up to the
irrelevant sign convention for translation.
A linear lift of this translation to $M$ has an $m$th tensor
power on $M^{\otimes m}$. Compare it, via $\psi$, with the
canonical deck lift on $[m]^*N$. The discrepancy is an
automorphism of a line bundle on the proper connected curve
$E$, and is therefore a scalar in $k^\times$.
Rescaling the lift removes this discrepancy. There are exactly
$m$ such rescalings. These compatible lifts form a group
\[
       1\longrightarrow\mu_m\longrightarrow G_m
         \longrightarrow E[m]\longrightarrow1.        \tag{D.EC.4}
\]
The kernel acts by fibre scalars and is central.
Off the zero section the action is free: a nonzero translation
has no fixed point on the base, while a nontrivial fibre scalar
has no fixed nonzero vector. Its order equals $\deg(f_m)$.
Thus (D.EC.3), off its zero section, is the corresponding
connected $G_m$-torsor.

Here is the exact commutator and exponent calculation, including
characteristic two. Apply the known theta-pairing theorem to
$M_0=\mathcal O_E(mO)$, whose degree $m$ is prime to $p$.
The translation kernel of $M_0$ is exactly $E[m]$.
Its commutator pairing is therefore a perfect alternating
pairing on $(\mathbb Z/m)^2$.
Since $M=M_0^{\otimes-3}$, tensoring lifts shows directly
that the commutator pairing for $M$ on $E[m]$ is the $(-3)$rd
power of this pairing. It is still perfect because $(3,m)=1$.
We have not applied a separable-type assertion to a line of
degree divisible by $p$ when $p=3$.

Give $N$ its symmetric inversion lift normalized to be the
identity on the fibre over $O$. Give $M=N^m$ the induced lift.
Both inversion structures in (D.EC.2) are normalized at $O$,
so $\psi$ respects them: their possible discrepancy is a
constant equal to one at $O$.
Conjugation by this inversion therefore defines an automorphism
$\iota$ of $G_m$ which induces $u\mapsto-u$ on $E[m]$ and
fixes $\mu_m$ pointwise. For $g\in G_m$ one has $g^m\in\mu_m$,
and $\iota(g)=g^{-1}c$ for some $c\in\mu_m$.
It follows that
\[
               g^m=\iota(g^m)=\iota(g)^m=g^{-m}.
                                                               \tag{D.EC.5}
\]
As $m$ is odd, $g^m=1$. This argument does not divide by two
in $k$ and does not invoke the characteristic restrictions
of the later symmetric-theta sections of Mumford's paper.

Choose the central generator $z$ to have the indicated
normal-parameter character $\zeta$.
Perfection of the commutator pairing permits a basis of
$E[m]$ whose lifts $x,y$ satisfy $[x,y]=z^3$.
The exponent calculation gives $x^m=y^m=1$.
The elements $x^ay^bz^c$, $0\le a,b,c<m$, are distinct,
since their images in $E[m]$ first determine $a,b$ and
then their central coordinate determines $c$.
They exhaust $G_m$. This proves (D.EC.1).

Finally we check that the torsor stays connected on the
fixed completed puncture. Compose $f_m$ with
$Y_0\to\operatorname{Spec}A$, where $A$ is the cone ring,
and take its Stein factorization. Its finite algebra
$A_m=\Gamma(\operatorname{Tot}(M),\mathcal O)$ is a normal
domain: global sections of an integral normal scheme are
integrally closed in their fraction field. Its field has
degree $m^3$ over $\operatorname{Frac}A$, since the Stein map
is an isomorphism generically. There is only one point of
$\operatorname{Spec}A_m$ over the vertex. Indeed its preimage
in $\operatorname{Tot}(M)$ is the connected zero section,
and the proper Stein map is surjective.
The puncture of this finite normalization is the torsor just
constructed, since $f_m$ is finite there. These uses of Stein
factorization are in its Noetherian form
\citep[Theorem~37.53.4, Tag~03H0]{D_StacksEllipticThetaCone}.

The completed algebra $A_m\otimes_A R$ is thus the completion
of one excellent normal local domain. By Tag~0C23 it is a
normal local domain, not a product. Its generic degree remains
$m^3$ by flat base change. Its puncture is the connected finite
\'etale $G_m$-torsor $U_m\to U$.
The resolution obtained by base change from
$\operatorname{Tot}(M)$ is regular by Tags~0BG4, and its
single exceptional curve and its map to $E$ are unchanged.
At the generic exceptional point, (D.EC.3) has ramification
index $m$ and residue degree $m^2$. Its inertia is exactly
the fibre-scalar subgroup $\langle z\rangle$.
\end{proof}

\begin{proposition}[Proved: unbounded actual unitary Brauer periods on one fixed puncture]
\label{D:fixed-elliptic-cone-unitary-geometric-periods}
For the one fixed complete surface $R$ of
Lemma~\ref{D:elliptic-cone-Heisenberg-covers}, there are
continuous finite-image representations
\[
 \rho_j:\pi_1(U)\longrightarrow\operatorname{PSU}_{m_j}(q_j)
             \longrightarrow\operatorname{PGL}_{m_j}(\mathbb F_{q_j^2})
                                                               \tag{D.EC.6}
\]
whose associated geometric Brauer classes on $U$ have
exact order $m_j$, where $m_j\to\infty$ at fixed $p$.
One can take, for $j\ge1$,
\[
 q_j=p^{2^j},\qquad
 m_j=\begin{cases}q_j+1,&p=2,\\ (q_j+1)/2,&p\ne2.\end{cases}
                                                               \tag{D.EC.7}
\]
Their images are the tame Heisenberg groups of order $m_j^3$
in (D.EC.1). In particular this is an assertion about actual
projective local systems and their geometric Brauer periods;
it is not just an assertion about an abstract scalar extension.
\end{proposition}

\begin{proof}
Fix $m,q$ from (D.EC.7). Then $m$ is odd, $m\ge5$,
$(m,6p)=1$, and $m\mid q+1$. For the assertion about three,
if $p=3$ then $q+1\equiv1\pmod3$; otherwise $q\equiv1\pmod3$.
For odd $p$, the even exponent also gives $q\equiv1\pmod8$,
so $(q+1)/2$ is odd. These observations prove all the stated
coprimality assertions. The sequence is plainly unbounded.

Put $K=\mathbb F_{q^2}\subset k$, and choose $\zeta\in K$
of exact order $m$. Its Hermitian norm is one.
On the standard orthonormal basis indexed by $j\in\mathbb Z/m$
define
\[
 Xe_j=e_{j+1},\qquad Ze_j=\zeta^j e_j,\qquad
 Ye_j=\zeta^{-3j(j-1)/2}e_j.                            \tag{D.EC.8}
\]
The quadratic exponent is well defined modulo $m$ since
$m$ is odd. All three matrices preserve the standard
nondegenerate Hermitian form. Their determinants are
\[
 \det X=1,\qquad \det Z=\zeta^{m(m-1)/2}=1,
 \qquad \det Y=\zeta^{-3m(m-1)(m-2)/6}=1.               \tag{D.EC.9}
\]
For the last equality, $(m,6)=1$ implies that
$(m-1)(m-2)/6$ is an integer.
Thus $X,Y,Z\in\operatorname{SU}_m(q)$, with the exact
unitary and projective conventions of
\citet[pp.~114--116, Lemma~10.1]{Taylor1992ClassicalB}.
Direct calculation gives
\[
 \begin{aligned}
 &X^m=Y^m=Z^m=I,\\
 &[X,Y]=\zeta^{-3}Z^3,\qquad [X,Z]=\zeta^{-1}I,
          \qquad [Y,Z]=I .
 \end{aligned}                                                \tag{D.EC.10}
\]
Hence $x\mapsto\overline X$, $y\mapsto\overline Y$,
$z\mapsto\overline Z$ defines a homomorphism
$G_m\to\operatorname{PSU}_m(q)$.
If $X^aY^bZ^c$ is scalar, its coordinate-line permutation
first gives $a=0$. The second finite difference, modulo $m$,
of its diagonal exponent $-3b j(j-1)/2+cj$ is $-3b$.
Thus $b=0$, since $(3,m)=1$, and then the first difference
gives $c=0$. The homomorphism is injective.
Composing it with the quotient supplied by the lemma gives
(D.EC.6). The stated $m,q$ lie in the usual nonabelian simple
unitary ranges; simplicity is not used in the argument.

Let $\alpha\in\operatorname{Br}(U)$ be the class of the
Azumaya algebra obtained from the $G_m$-torsor by descent
of $\operatorname{Mat}_m(\mathcal O_{U_m})$ under these
projective matrices. The section
$x^ay^bz^c\mapsto X^aY^bZ^c$, with exponents between zero
and $m-1$, has scalar multiplication defects in $\mu_m$,
by (D.EC.10). Consequently its obstruction on $U$ is the
image of a class in $H^2(U,\mu_m)$, and
\[
                              m\alpha=0.               \tag{D.EC.11}
\]

We prove the opposite inequality in the \emph{geometric}
Brauer group. Let $H=\langle x,z\rangle\simeq(\mathbb Z/m)^2$.
Restrict the torsor to its intermediate quotient by $H$,
and complete the exceptional discrete valuation.
The resulting Galois field extension $\mathcal L/\mathcal K$
has group $H$, inertia $\langle z\rangle$, and residue
extension cyclic of degree $m$ with generator $x$.
Here the unique exceptional curve upstairs ensures that
the decomposition group is all of $H$.
Thus $\mathcal L^{\langle z\rangle}/\mathcal K$ is an
unramified extension of degree $m$, and
$\mathcal L^{\langle x\rangle}/\mathcal K$ is a totally
ramified cyclic extension of degree $m$.

Choose a Kummer generator $\xi$ of the unramified extension
with $x(\xi)=\zeta\xi$. Multiplying it by a power of a
uniformizer of $\mathcal K$ makes it a unit. Its residue
still generates the degree-$m$ residue extension: reduction
preserves the primitive eigencharacter, and $m$ is invertible.
Write $a=\xi^m\in\mathcal K^\times$.
Choose a uniformizer $s_0$ of the totally ramified field.
The action of $z$ on its normal line is $\zeta$, by our
choice in the lemma. The eigenspace projection
\[
       s=\frac1m\sum_{i=0}^{m-1}\zeta^{-i}z^i(s_0)
                                                               \tag{D.EC.12}
\]
is a uniformizer, satisfies $z(s)=\zeta s$, and is fixed
by $x$. Therefore $b=s^m\in\mathcal K^\times$ has
$v_{\mathcal K}(b)=1$.

The actual descended algebra can now be identified without
any inference from the order of an abstract group-cohomology
class. On $\mathcal L\otimes_k\operatorname{Mat}_m(k)$ the
semilinear descent actions of $x,z$ are their actions on
$\mathcal L$ followed by conjugation by $X,Z$, respectively.
The two matrices
\[
             A=\xi Z,\qquad B=s^{-1}X
                                                               \tag{D.EC.13}
\]
are invariant under both actions, since
$XZX^{-1}=\zeta^{-1}Z$ and $ZXZ^{-1}=\zeta X$.
They satisfy
\[
             A^m=aI,\qquad B^m=b^{-1}I,
                       \qquad AB=\zeta BA.             \tag{D.EC.14}
\]
The $m^2$ matrices $Z^iX^j$, $0\le i,j<m$, are a basis
of the matrix algebra over $\mathcal L$: the shifts give
the different diagonals and the primitive-root Vandermonde
matrix separates their diagonal coefficients.
Consequently $A^iB^j$ are linearly independent and span
the descended algebra over $\mathcal K$.
It is precisely the symbol algebra $(a,b^{-1})_\zeta$.

We use the known cyclic-algebra norm criterion, with no
primality restriction on $m$:
\citet[Proposition~4.7.3 and Corollaries~4.7.4--4.7.5,
pp.~109--111]{D_GilleSzamuely2006Cone}.
It identifies the class of this algebra, up to the harmless
inverse choice of cyclic generator, with the class of
$b^{-1}$ modulo norms from the unramified degree-$m$
extension $\mathcal K(\xi)/\mathcal K$.
Its $e$th multiple is therefore trivial only if $b^{-e}$
is such a norm. Every norm from this unramified extension
has valuation divisible by $m$: for $u\ne0$, all $m$
conjugates have the same integral valuation, and the
valuation of their product is $m v(u)$.
Since $v(b)=1$, this forces $m\mid e$.
Together with (D.EC.11), it proves that the restriction of
$\alpha$ to this actual valued-field extension, and hence
$\alpha$ itself on $U$, has exact order $m$.

In particular Proposition~\ref{D:single-curve-Brauer-residue-cover}
identifies its residue with a class of exact order $m$ in
$H^1(E,\mathbb Z/m)$. This last observation is a consequence
of the established geometric order, not an appeal to the
prime-order-only eigensheaf assertions in the literature.
The fixed germ $R$ and fixed exceptional curve $E$ have
not varied with $m$ or $q$. This completes the proof.
\end{proof}

\paragraph{Scope of the diagnostic.}
The image of (D.EC.6) is the explicitly displayed nilpotent
group $G_m$, not the whole simple unitary group. Thus this
does not construct epimorphisms onto the ambient simple
groups, does not contradict any quadratic representation
bound, and does not close an unbounded-period hereditary
count. It proves that a uniform bound on geometric Brauer
period cannot be inferred merely from a unitary form,
even for actual tame projective local systems on one fixed
complete simple elliptic surface puncture.
